\documentclass[11pt]{amsart}

\usepackage[margin=1.1in]{geometry}
\usepackage{amsmath,amssymb,amsthm}
\usepackage{enumitem}
\usepackage{graphicx}
\usepackage[colorlinks=true,linkcolor=blue,citecolor=blue,urlcolor=blue]{hyperref}

\newtheorem{theorem}{Theorem}
\newtheorem{lemma}{Lemma}
\newtheorem{corollary}{Corollary}

\newcommand{\lip}{\operatorname{lip}}
\newcommand{\Lip}{\operatorname{Lip}}
\newcommand{\eps}{\varepsilon}
\newcommand{\R}{\mathbb{R}}
\newcommand{\HH}{\mathcal{H}}

\title[Metric wedges and locally flat separation]{Metric wedges of compact purely 1-unrectifiable spaces have locally flat uniform separation}
\author{arxiv\_automatic\_math run fa\_banach\_001}
\date{2026-06-28}

\begin{document}
\maketitle

\section*{Status}

Result type: partial result.  Status: likely valid pending human review.

This packet addresses Question~2.7 (TeX label \texttt{q:spu p1u}) in
Aliaga's exposition \cite{Aliaga2606}.  The question asks whether, for a
purely $1$-unrectifiable metric space $M$, the algebra $\lip(M)$ of locally
flat Lipschitz functions separates points of $M$ uniformly.  The packet proves
a positive answer for metric wedges of compact purely $1$-unrectifiable
pointed spaces.  This is not a full solution of the source question.

\section*{Open-problem evidence}

The source question appears on page 6 of the rendered source PDF:

\begin{center}
\includegraphics[width=\textwidth]{figures/open_problem_crop.png}
\end{center}

Transcription: Let $M$ be a purely $1$-unrectifiable metric space. Does
$\lip(M)$ separate points of $M$ uniformly?

\section*{Definitions}

A real-valued Lipschitz function $f$ on a metric space $M$ is locally flat if
for every $p\in M$,
\[
  \lim_{r\downarrow 0}
  \sup\left\{
    \frac{|f(u)-f(v)|}{d(u,v)}:
    0<d(u,v),\ d(u,p)<r,\ d(v,p)<r
  \right\}=0.
\]
The space $\lip(M)$ consists of locally flat Lipschitz functions.
It separates points uniformly if for every $x,y\in M$ and every $\eps>0$ there
is $f\in\lip(M)$ with $\Lip(f)\leq 1$ and
\[
  f(y)-f(x)>d(x,y)-\eps.
\]

Let $\{(K_i,o_i)\}_{i\in I}$ be pointed metric spaces.  Their metric wedge
$W=\bigvee_{i\in I} K_i$ is obtained by identifying all basepoints $o_i$ to a
single point $o$.  If $a,b$ lie in the same branch $K_i$, their distance in
$W$ is $d_i(a,b)$.  If $a\in K_i$, $b\in K_j$, $i\neq j$, then
\[
  d_W(a,b)=d_i(a,o_i)+d_j(b,o_j).
\]

\section*{Proof intuition}

For a compact purely $1$-unrectifiable branch, the AGPP theorem
\cite{AGPP} gives a locally flat, $1$-Lipschitz function that almost realizes
the distance between any two points.  If such a branch function is normalized
to vanish at the wedge basepoint, extending it by zero to all other branches
stays $1$-Lipschitz because the wedge distance to any other branch includes
the whole distance back to the basepoint.  For points in two different
branches, use one normalized separator on each branch, with opposite signs.
The only point at which local flatness could fail is the common basepoint; it
does not fail because each branch function is $o(d(\cdot,o))$ near the base,
so for two nearby points in distinct branches the numerator is bounded by the
sum of two small radial errors.

The original question remains open outside this wedge class.

\section*{Main theorem}

\begin{theorem}
\label{thm:wedge-spu}
Let $\{(K_i,o_i)\}_{i\in I}$ be any family of pointed compact metric spaces.
Assume every $K_i$ is purely $1$-unrectifiable.  Then
\[
  W=\bigvee_{i\in I} K_i
\]
has the following property: $\lip(W)$ separates points of $W$ uniformly.
\end{theorem}

\begin{proof}
We use the compact theorem of Aliaga, Gartland, Petitjean, and Prochazka:
if $K$ is compact and purely $1$-unrectifiable, then $\lip(K)$ separates
points of $K$ uniformly \cite[Theorem 2.10]{Aliaga2606},
\cite[Theorem A]{AGPP}.

Fix $x,y\in W$ and $\eps>0$.  If $x=y$, the zero function works, so assume
$x\neq y$.

First suppose that $x,y$ belong to the same branch $K_i$, where the basepoint
$o$ is identified with $o_i$ if one of the points is the wedge point.  By the
compact theorem there is $h\in\lip(K_i)$ such that $\Lip(h)\leq 1$ and
\[
  h(y)-h(x)>d_i(x,y)-\eps .
\]
Replacing $h$ by $h-h(o_i)$, assume $h(o_i)=0$.  Define $F:W\to\R$ by
$F=h$ on $K_i$ and $F=0$ on all other branches.  The function is
$1$-Lipschitz: inside $K_i$ this is clear; if $a\in K_i$ and $b$ is in another
branch, then
\[
  |F(a)-F(b)|=|h(a)|\leq d_i(a,o_i)\leq d_W(a,b);
\]
and if both points are outside $K_i$, the values are both zero.  Thus
$F(y)-F(x)>d_W(x,y)-\eps$.

It remains in this case only to verify local flatness.  Away from the wedge
point, a small ball is contained in a single branch, so local flatness follows
from that of $h$ or from the fact that $F$ is locally constant.  At the wedge
point, let $\eta>0$.  Since $h\in\lip(K_i)$ and $h(o_i)=0$, there is
$\delta>0$ such that $|h(a)|\leq\eta d_i(a,o_i)$ whenever
$a\in K_i$ and $d_i(a,o_i)<\delta$, and also such that
$|h(a)-h(b)|\leq\eta d_i(a,b)$ for $a,b\in K_i$ within $\delta$ of $o_i$.
For any $u,v\in W$ within $\delta$ of $o$, if both are in $K_i$ the latter
estimate applies; if exactly one is in $K_i$, say $u\in K_i$, then
\[
  |F(u)-F(v)|=|h(u)|\leq \eta d_i(u,o_i)\leq \eta d_W(u,v);
\]
and if neither is in $K_i$, the numerator is zero.  Hence $F$ is locally flat
at $o$.

Now suppose $x\in K_i\setminus\{o\}$ and $y\in K_j\setminus\{o\}$ with
$i\neq j$.  Put $r_x=d_i(x,o_i)$ and $r_y=d_j(y,o_j)$.  By the compact theorem
choose $h_i\in\lip(K_i)$ and $h_j\in\lip(K_j)$ with Lipschitz constants at
most $1$, normalized by $h_i(o_i)=h_j(o_j)=0$, such that
\[
  h_i(x)>r_x-\eps/2,
  \qquad
  h_j(y)>r_y-\eps/2 .
\]
Define $F:W\to\R$ by $F=-h_i$ on $K_i$, $F=h_j$ on $K_j$, and $F=0$ on all
other branches.  Then
\[
  F(y)-F(x)=h_j(y)+h_i(x)>r_y+r_x-\eps=d_W(x,y)-\eps .
\]

The same radial estimates prove that $F$ is $1$-Lipschitz.  Inside a selected
branch this is the Lipschitz bound for $h_i$ or $h_j$.  Between the two
selected branches,
\[
  |F(a)-F(b)|
  \leq |h_i(a)|+|h_j(b)|
  \leq d_i(a,o_i)+d_j(b,o_j)=d_W(a,b),
\]
and the estimates against an unselected branch are the same with one term
missing.  Local flatness away from $o$ is branch-local.  At $o$, choose
$\delta$ so that both selected functions have the corresponding radial and
same-branch local-flat estimates with error $\eta$.  If two nearby points lie
in different selected branches, the preceding display improves to
$|F(u)-F(v)|\leq \eta(d_W(u,o)+d_W(v,o))=\eta d_W(u,v)$; all other cases are
easier.  Hence $F\in\lip(W)$.

This proves uniform separation for all pairs of points in $W$.
\end{proof}

\begin{lemma}
\label{lem:complete}
If each branch $K_i$ is complete, then the metric wedge $W$ is complete.  In
particular, the wedges in Theorem~\ref{thm:wedge-spu} are complete.
\end{lemma}

\begin{proof}
Let $(z_n)$ be a Cauchy sequence in $W$ and put $r_n=d_W(z_n,o)$.  If
$r_n\to 0$, then $z_n\to o$.  Otherwise some tail has points with radius
bounded below.  Choose $\alpha>0$ and infinitely many $n$ with $r_n>2\alpha$.
Since $(z_n)$ is Cauchy, there is $N$ such that $d_W(z_m,z_n)<\alpha$ for
$m,n\geq N$.  Taking one $n\geq N$ with $r_n>2\alpha$, every later point must
lie in the same branch as $z_n$; a point in another branch, including the base,
would be at distance at least $r_n>\alpha$ from $z_n$.  The tail is therefore
Cauchy in one complete branch and converges there.
\end{proof}

\begin{corollary}
The source question has an affirmative answer for every complete metric wedge
of compact purely $1$-unrectifiable pointed metric spaces.  Such wedges may be
noncompact, for instance when there are infinitely many nontrivial branches
with diameters bounded away from zero.
\end{corollary}

\begin{proof}
The affirmative answer is Theorem~\ref{thm:wedge-spu}, and completeness is
Lemma~\ref{lem:complete}.  Noncompactness follows, for example, by taking
infinitely many two-point branches of the same positive length.
\end{proof}

\section*{Verification notes}

The proof has no computational dependency.  The key checks are:
\begin{itemize}[leftmargin=*]
\item Normalizing each branch function by $h(o_i)=0$ gives
  $|h(a)|\leq d_i(a,o_i)$, so extension by zero is $1$-Lipschitz across
  branches.
\item In the two-branch case the signs make
  $F(y)-F(x)$ add the two radial branch distances, exactly the wedge distance.
\item Local flatness at the wedge point follows from radial smallness
  $|h(a)|=o(d_i(a,o_i))$ for each selected branch.
\end{itemize}

\section*{Novelty and limitations}

The local run indexes were searched for \texttt{2606.02918}, \texttt{q:spu},
\texttt{locally flat}, \texttt{separates points uniformly}, \texttt{metric
wedge}, and \texttt{purely 1-unrectifiable}.  No duplicate packet for this
question or wedge closure result was found.  A bounded web search on
2026-06-28 for exact phrases including ``lip(M) separates points'',
``locally flat Lipschitz functions'' with ``metric wedge'', and ``purely
1-unrectifiable'' with ``wedge'' found the AGPP compact theorem but no exact
wedge result.

This packet does not solve Question~2.7 for arbitrary purely
$1$-unrectifiable metric spaces.  It proves a closure theorem that manufactures
possibly noncompact positive examples from compact positive ones.

\section*{Human review recommendation}

Please check the basepoint local-flatness argument and the completeness lemma.
The rest of the proof is a direct branchwise use of the compact theorem and
the defining formula for the wedge metric.

\begin{thebibliography}{9}

\bibitem{Aliaga2606}
R. J. Aliaga,
\emph{Lipschitz-free spaces and purely 1-unrectifiable metric spaces},
arXiv:2606.02918, 2026.
\url{https://arxiv.org/abs/2606.02918}

\bibitem{AGPP}
R. J. Aliaga, C. Gartland, C. Petitjean, and A. Prochazka,
\emph{Purely 1-unrectifiable metric spaces and locally flat Lipschitz functions},
Trans. Amer. Math. Soc. \textbf{375} (2022), 3529--3567.
\url{https://arxiv.org/abs/2103.09370}

\end{thebibliography}

\end{document}
